\documentclass[11pt]{article}

% --- Paquetes estándar (compila con pdflatex en cualquier TeX Live) ---
\usepackage[utf8]{inputenc}
\usepackage[T1]{fontenc}
\usepackage[margin=2.5cm]{geometry}
\usepackage{amsmath,amssymb}
\usepackage{booktabs}
\usepackage{array}
\usepackage{float}
\usepackage{xcolor}
\usepackage[colorlinks=true,linkcolor=blue,urlcolor=blue]{hyperref}

\newcommand{\best}[1]{\textbf{#1}}

\title{\textbf{Mejora del desempeño en datos reales (Dataset híbrido) \\
para segmentación 3D de tumores en DBT}\\[4pt]
\large Fine-tuning sim$\rightarrow$real, ensembles y análisis del límite de datos}
\author{Gabriel Guerra}
\date{\today}

\begin{document}
\maketitle

\begin{abstract}
\noindent Se documenta una campaña de experimentos orientada a mejorar el desempeño
de segmentación de tumores en imágenes \emph{reales} de tomosíntesis mamaria (DBT),
correspondientes al dataset híbrido \textbf{D004} (\texttt{Dataset\_Both\_RealWorld}),
cuyo conjunto de test está compuesto al 100\% por casos reales y que partía de un
Dice de $\approx 0.35$. Mediante una metodología de \emph{fine-tuning} de dominio
sintético$\rightarrow$real, función de pérdida orientada a \emph{recall} (Tversky),
validación cruzada de 5 folds y \emph{ensembles} de arquitecturas, se elevó el Dice
de test a \best{0.518} (mediana 0.594), eliminando todos los casos con detección nula
(de 7/10 a 0/10). Se analiza el techo de desempeño alcanzable: con solo 20 casos reales
disponibles, el límite oráculo es $\approx 0.53$, por lo que objetivos como $0.6$
no son alcanzables sin incorporar más datos reales. Todos los experimentos son
reproducibles y se reportan también los enfoques que \emph{no} mejoraron.
\end{abstract}

\section{Problema y diagnóstico}
El proyecto compara cuatro arquitecturas sobre cuatro datasets de DBT. El dataset
\textbf{D004} es el único cuya evaluación de test es 100\% \emph{out-of-distribution}
(casos reales \texttt{real\_dbt\_*}), mientras que el entrenamiento es mayoritariamente
sintético. El bajo desempeño inicial ($\sim$0.35) se diagnosticó como un
\textbf{gap de dominio sintético$\rightarrow$real}, y no como un problema de
arquitectura ni de geometría:
\begin{itemize}\itemsep2pt
  \item Todas las imágenes (sintéticas y reales) comparten dimensiones idénticas
        ($512\times256\times15$, \emph{spacing} isotrópico $1$\,mm); el plan de nnU-Net
        no aplica \emph{resampling}. Se descartó cualquier artefacto geométrico.
  \item La distribución de intensidades difiere: los sintéticos presentan una cola
        brillante (máx.\ $\sim$4--5$\sigma$) ausente en los reales (máx.\ $\sim$2$\sigma$).
  \item El fallo es de \emph{recall}: el modelo sub-segmenta, no genera falsos positivos.
\end{itemize}

\section{Datos}
\begin{itemize}\itemsep2pt
  \item \textbf{Sintéticos:} 120 casos \texttt{dbt\_001..120}.
  \item \textbf{Reales:} 20 casos \texttt{real\_dbt\_001..020} (\emph{único} conjunto real del proyecto; 20 pacientes distintos).
  \item \textbf{Partición D004:} 10 reales de desarrollo (\texttt{imagesTr}) / 10 reales de test (\texttt{imagesTs}, \emph{held-out}).
  \item \textbf{Sin fuga de datos (leakage):} pacientes distintos y sintéticos de fuente independiente (confirmado), por lo que el modelo base sintético no contiene información de los pacientes reales de test.
\end{itemize}

\section{Metodología}
\subsection{Fine-tuning sim$\rightarrow$real}
Se parte de un modelo entrenado en datos sintéticos (\emph{warm-start}) y se ajusta
sobre los 10 casos reales de desarrollo con \emph{learning rate} reducido. La pérdida
combina entropía cruzada con \textbf{Tversky} ($\alpha{=}0.3,\ \beta{=}0.7$), que
penaliza con más fuerza los falsos negativos para incrementar el \emph{recall}:
\begin{equation}
\mathrm{TI} = \frac{\mathrm{TP}}{\mathrm{TP} + \alpha\,\mathrm{FP} + \beta\,\mathrm{FN}},
\qquad \mathcal{L} = \mathrm{CE} + (1-\mathrm{TI}).
\end{equation}

\subsection{Validación cruzada y selección sin leakage}
Sobre los 10 casos de desarrollo se usa una partición de \textbf{5 folds}. Las
predicciones \emph{out-of-fold} (OOF) dan una señal de validación limpia con la que se
seleccionan hiperparámetros del ensemble (pesos, umbral, post-proceso) \emph{sin tocar
el test}. El conjunto de test (10 reales) se evalúa una sola vez con la configuración
ya fijada.

\subsection{Ensemble}
Se promedian las probabilidades de varias redes (cada una con su ensemble de 5 folds).
Los pesos por red y el umbral se eligen maximizando el Dice sobre el OOF de desarrollo.

\section{Resultados}
\subsection{Progresión del desempeño en test (10 reales)}
\begin{table}[H]\centering
\begin{tabular}{lccc}
\toprule
Etapa & Dice (media) & Mediana & Ceros \\
\midrule
Zero-shot (sintético sobre real)        & 0.163 & 0.000 & 7/10 \\
Mezcla naive (histórico D004)            & 0.353 & ---   & 5/10 \\
nnU-Net fine-tuning (Dice+CE)            & 0.403 & 0.387 & 2/10 \\
nnU-Net + Tversky (1 fold)               & 0.465 & 0.524 & 2/10 \\
nnU-Net + Tversky (ensemble 5 folds)     & 0.482 & 0.636 & 2/10 \\
nnU-Net + umbral/combinación             & 0.505 & 0.556 & 1/10 \\
\best{Ensemble 3 redes (nn+att+bce)}     & \best{0.518} & \best{0.594} & \best{0/10} \\
\bottomrule
\end{tabular}
\caption{Evolución del Dice de test. El mejor resultado es el ensemble de las tres
redes más fuertes, que además detecta el 100\% de las lesiones (0 casos nulos).}
\end{table}

\subsection{Fine-tuning por arquitectura (single model)}
\begin{table}[H]\centering
\begin{tabular}{lccc}
\toprule
Arquitectura & Dice (media) & Mediana & Ceros \\
\midrule
nnU-Net (Tversky, 5-fold + combinación) & 0.505 & 0.556 & 1/10 \\
DynUNet (MONAI)                          & 0.475 & ---   & 1/10 \\
Attention U-Net 3D                       & 0.407 & 0.425 & 0/10 \\
SegResNet (MONAI)                        & 0.395 & ---   & 1/10 \\
U-Net (BCE$\rightarrow$Tversky)          & 0.394 & 0.449 & 1/10 \\
V-Net (MONAI)                            & 0.251 & ---   & 4/10 \\
\bottomrule
\end{tabular}
\caption{Desempeño individual tras fine-tuning. Las redes son complementarias:
nnU-Net domina los casos fáciles; U-Net (BCE) y Attention rescatan casos de bajo
contraste que las demás fallan.}
\end{table}

\subsection{Ensembles}
\begin{table}[H]\centering
\begin{tabular}{lccc}
\toprule
Ensemble (pesos elegidos en OOF dev) & Dice (media) & Mediana & Ceros \\
\midrule
\best{3 redes: nn$\cdot$2 + att$\cdot$2 + bce$\cdot$3} & \best{0.518} & 0.594 & 0/10 \\
4 redes (+ DynUNet)                                    & 0.508 & 0.581 & 0/10 \\
6 redes (+ SegResNet, V-Net)                           & 0.508 & 0.581 & 0/10 \\
nn + DynUNet                                           & 0.500 & 0.616 & 2/10 \\
\bottomrule
\end{tabular}
\caption{``Menos es más'': el ensemble de las tres redes complementarias supera a
configuraciones con más arquitecturas. SegResNet y V-Net resultan redundantes/débiles
y la selección en dev les asigna peso nulo.}
\end{table}

\subsection{Resultado final por caso}
\begin{table}[H]\centering
\begin{tabular}{lc|lc}
\toprule
Caso & Dice & Caso & Dice \\
\midrule
real\_dbt\_005 & 0.778 & real\_dbt\_013 & 0.590 \\
real\_dbt\_006 & 0.652 & real\_dbt\_014 & 0.519 \\
real\_dbt\_007 & 0.461 & real\_dbt\_015 & 0.727 \\
real\_dbt\_010 & 0.746 & real\_dbt\_016 & 0.598 \\
real\_dbt\_019 & 0.050 & real\_dbt\_020 & 0.064 \\
\bottomrule
\end{tabular}
\caption{Dice por caso del ensemble final de 3 redes (media 0.518, mediana 0.594).
Los casos 019 y 020 (tumores de muy bajo contraste) son el cuello de botella, pero
ya no quedan en cero.}
\end{table}

\section{Enfoques evaluados que no mejoraron}
Por transparencia, se reportan los experimentos que \emph{no} aportaron:
\begin{itemize}\itemsep2pt
  \item \textbf{Backbone 2.5D pre-entrenado en ImageNet} (U-Net con encoder ResNet34,
        slices vecinos como canales): Dice 0.280. La transferencia sintético$\rightarrow$real
        fue nula y se pierde el contexto inter-slice del 3D.
  \item \textbf{Mezcla del 2.5D con nnU-Net:} cualquier peso del modelo 2.5D \emph{baja}
        la media del ensemble (de 0.505 a 0.459); un miembro débil perjudica.
  \item \textbf{CLAHE (realce de contraste):} aunque mejora la separabilidad
        tumor-fondo (Cohen's $d$ del caso 020: $0.79\rightarrow1.28$), no se tradujo en
        Dice (U-Net 0.395, Attention 0.434); el realce global introduce ruido.
  \item \textbf{Arquitecturas adicionales} (SegResNet, V-Net): demasiado débiles con
        tan pocos datos; no aportan al ensemble.
\end{itemize}

\section{Análisis del límite de desempeño}
Se calculó el \textbf{límite oráculo}: tomando, para cada caso, el mejor de todos los
modelos entrenados, la media máxima alcanzable es $\approx 0.53$. Es decir, ni con un
selector perfecto (no realizable sin incurrir en leakage) se superaría ese valor.
La causa es estructural:
\begin{itemize}\itemsep2pt
  \item Solo existen \textbf{20 casos reales} en todo el proyecto.
  \item Los casos 019 y 020 son tumores reales de muy bajo contraste; el 020 no era
        detectado por \emph{ningún} modelo individual (miss genuino, no error de etiqueta).
\end{itemize}
Por lo tanto, \textbf{objetivos como $0.6$ no son alcanzables por vía algorítmica}
con los datos disponibles; el factor limitante es la cantidad de datos reales.

\section{Conclusiones}
\begin{itemize}\itemsep2pt
  \item El fine-tuning sim$\rightarrow$real con pérdida Tversky y \emph{ensemble} de las
        tres redes más fuertes mejora el Dice de test de $0.35$ a \best{0.518}
        ($+0.17$), eliminando todos los casos con detección nula.
  \item La metodología es \emph{limpia} (selección de hiperparámetros en OOF de
        desarrollo, sin leakage) y reproducible.
  \item Se demostró un techo de datos ($\approx 0.53$): la contribución incluye tanto
        la mejora metodológica como el análisis honesto del límite.
\end{itemize}

\section{Trabajo futuro}
\begin{itemize}\itemsep2pt
  \item \textbf{Más datos reales} (única vía real a $>0.53$), idealmente de bajo contraste.
  \item \textbf{Semi-supervisado / pseudo-labeling} si existieran casos reales sin etiquetar.
  \item \textbf{Aumentación tipo copy-paste de tumores} para sintetizar variedad real.
  \item \textbf{Evaluación con LOOCV} sobre los 20 reales y \textbf{métrica de detección}
        (sensibilidad por lesión), clínicamente más relevante; el ensemble actual ya
        detecta el 100\% de las lesiones.
\end{itemize}

\appendix
\section{Reproducibilidad}
Scripts principales (bajo \texttt{experiments/}):
\begin{itemize}\itemsep1pt
  \item nnU-Net fine-tuning: \texttt{nnunet\_official/scripts/trainers/nnUNetTrainer\_FT.py},
        \texttt{nnUNetTrainer\_FT\_Tversky.py}.
  \item Redes custom (Attention, U-Net): \texttt{custom\_nets\_ft/full\_pipeline.py}.
  \item Arquitecturas adicionales + ensemble: \texttt{ensemble\_plus/run\_all.py},
        \texttt{best\_ensemble.py}.
  \item Backbone 2.5D: \texttt{backbone\_2p5d/train\_eval.py}.
  \item CLAHE: \texttt{clahe/clahe\_pipeline.py}.
\end{itemize}
Resultados en formato JSON junto a cada script (\texttt{*\_result.json}, \texttt{FINAL.json},
\texttt{best\_ensemble.json}).

\end{document}
